
% ----------------------------------------------------------------------
%                   LATEX TEMPLATE FOR PhD THESIS
% ----------------------------------------------------------------------

% based on Harish Bhanderi's PhD/MPhil template, then Uni Cambridge
% http://www-h.eng.cam.ac.uk/help/tpl/textprocessing/ThesisStyle/
% corrected and extended in 2007 by Jakob Suckale, then MPI-CBG PhD programme
% and made available through OpenWetWare.org - the free biology wiki


%: Style file for Latex
% Most style definitions are in the external file PhDthesisPSnPDF.
% In this template package, it can be found in ./Latex/Classes/
\documentclass[twoside,11pt]{Latex/Classes/PhDthesisPSnPDF}

\makeatletter
        \setlength{\@fptop}{0pt}
\makeatother


%: Macro file for Latex
% Macros help you summarise frequently repeated Latex commands.
% Here, they are placed in an external file /Latex/Macros/MacroFile1.tex
% An macro that you may use frequently is the figuremacro (see introduction.tex)
% This file contains macros that can be called up from connected TeX files
% It helps to summarise repeated code, e.g. figure insertion (see below).

% insert a centered figure with caption and description
% parameters 1:filename, 2:title, 3:description and label
\newcommand{\figuremacro}[3]{
	\begin{figure}[htbp]
		\centering
		\includegraphics[width=1\textwidth]{#1}
		\caption[#2]{\textbf{#2} - #3}
		\label{#1}
	\end{figure}
}

% insert a centered figure with caption and description AND WIDTH
% parameters 1:filename, 2:title, 3:description and label, 4: textwidth
% textwidth 1 means as text, 0.5 means half the width of the text
\newcommand{\figuremacroW}[4]{
	\begin{figure}[htbp]
		\centering
		\includegraphics[width=#4\textwidth]{#1}
		\caption[#2]{\textbf{#2} - #3}
		\label{#1}
	\end{figure}
}

% inserts a figure with wrapped around text; only suitable for NARROW figs
% o is for outside on a double paged document; others: l, r, i(inside)
% text and figure will each be half of the document width
% note: long captions often crash with adjacent content; take care
% in general: above 2 macro produce more reliable layout
\newcommand{\figuremacroN}[3]{
	\begin{wrapfigure}{o}{0.5\textwidth}
		\centering
		\includegraphics[width=0.48\textwidth]{#1}
		\caption[#2]{{\small\textbf{#2} - #3}}
		\label{#1}
	\end{wrapfigure}
}

% predefined commands by Harish
\newcommand{\PdfPsText}[2]{
  \ifpdf
     #1
  \else
     #2
  \fi
}

\newcommand{\IncludeGraphicsH}[3]{
  \PdfPsText{\includegraphics[height=#2]{#1}}{\includegraphics[bb = #3, height=#2]{#1}}
}

\newcommand{\IncludeGraphicsW}[3]{
  \PdfPsText{\includegraphics[width=#2]{#1}}{\includegraphics[bb = #3, width=#2]{#1}}
}

\newcommand{\InsertFig}[3]{
  \begin{figure}[!htbp]
    \begin{center}
      \leavevmode
      #1
      \caption{#2}
      \label{#3}
    \end{center}
  \end{figure}
}


%%% Local Variables: 
%%% mode: latex
%%% TeX-master: "~/Documents/LaTeX/CUEDThesisPSnPDF/thesis"
%%% End: 
\newcommand{\norm}[1]{\left\lVert#1\right\rVert}

\newcommand\numberthis{\addtocounter{equation}{1}\tag{\theequation}}

\newcommand{\eg}{\textit{e.g. }}
\newcommand{\ie}{\textit{i.e. }}

\usepackage[table]{xcolor}
%% The amssymb package provides various useful mathematical symbols
\usepackage{amssymb}

%% The amsthm package provides extended theorem environments
\usepackage{amsthm}

%% amsmath for math environment
\usepackage{amsmath}

\DeclareMathOperator*{\argmin}{arg\,min}
\DeclareMathOperator*{\argmax}{arg\,max}
\DeclareMathOperator*{\sign}{sign}
\DeclareMathOperator*{\infspie}{inf}


% to break equation
%\usepackage{mathpazo}
%\usepackage{mathptmx}
%\usepackage[mathpazo]{flexisym}
%\usepackage{breqn}

%% For clever reference
%\usepackage{cleveref}

%% color package
\usepackage{color}

%% figure package
\usepackage{epsf,graphicx}
\usepackage{epstopdf}
\usepackage{subfigure}	
\usepackage{transparent}

%% New environment to have some indent inside enumerate environment
\usepackage{enumitem}

%% To create acronym for proper glossary
\usepackage{acro}

%% To number the line in the article
\usepackage{lineno}

%% Environment to include table with notes
\usepackage{array}
\usepackage{threeparttable}
\usepackage{booktabs}
\usepackage{multirow}
\usepackage{siunitx}
\usepackage{hhline}

%% In order to change size of margin
\usepackage{geometry}
\usepackage{changepage}
\usepackage{lscape}
%% Colorpackage for table
\usepackage{colortbl}
\usepackage{tabularx}
\usepackage{arydshln}

%% To use URL referencing
\usepackage{url}
%\usepackage[hidelinks]{hyperref}

%% In order to draw some graphs
\usepackage{tikz,xifthen}
\usepackage{tikz-qtree}
\usetikzlibrary{decorations.pathmorphing} % noisy shapes
\usetikzlibrary{fit}					% fitting shapes to coordinates
\usetikzlibrary{backgrounds}	% drawing the background after the foreground
\usetikzlibrary{shapes,arrows,shadows}
\usetikzlibrary{calc,decorations.pathreplacing,decorations.markings,positioning}
\usetikzlibrary{snakes,decorations.text,shapes,patterns}
%\usepackage{scalefnt,lmodern,booktabs}

%% Paxkage for cross and tick symbols
\usepackage{pifont}
\newcommand{\cmark}{\color{green!60!black!80}\ding{51}}
\newcommand{\mmark}{{\color{green!60!black!80}\ding{51}}$^{!}$}
\newcommand{\xmark}{\color{red!60!black!80}\ding{55}}
\newcommand{\cmarksmall}{\color{green!60!black!80}\ding{51}}
\newcommand{\mmarksmall}{{\color{green!60!black!80}\ding{51}}$^{!}$}
\newcommand{\xmarksmall}{\color{red!60!black!80}\ding{55}}
\newcommand{\Conv}{\mathop{\scalebox{1.5}{\raisebox{-0.2ex}{$\ast$}}}}%

\definecolor{autoGuided}{rgb}{ 0.3765    0.7294    0.9412}
\newcommand{\autoGuidedColor}{(light-Blue)}
\definecolor{fullyAuto}{rgb}{ 0.0941    0.3843    0.6627}
\newcommand{\fullyAutoColor}{(dark-blue)}
\definecolor{semiAuto}{rgb}{ 0.0784    0.5059    0.1686}
\newcommand{\semiAutoColor}{(light-green)}
\definecolor{fullyGuided}{rgb}{ 0.4275    0.6902    0.3176}
\newcommand{\fullyGuidedColor}{(dark-green)}

\DeclareSIUnit\ppm{ppm}
\DeclareSIUnit\px{px}

\usepackage{ltxtable}
\usepackage{listings}
\usepackage[toc]{appendix}
 
\definecolor{codegreen}{rgb}{0,0.6,0}
\definecolor{codegray}{rgb}{0.5,0.5,0.5}
\definecolor{codepurple}{rgb}{0.58,0,0.82}
\definecolor{backcolour}{rgb}{0.95,0.95,0.92}
 
\lstdefinestyle{mystyle}{
    backgroundcolor=\color{backcolour},   
    commentstyle=\color{codegreen},
    keywordstyle=\color{magenta},
    numberstyle=\tiny\color{codegray},
    stringstyle=\color{codepurple},
    basicstyle=\footnotesize,
    breakatwhitespace=false,         
    breaklines=true,                 
    captionpos=b,                    
    keepspaces=true,                 
    numbers=left,                    
    numbersep=5pt,                  
    showspaces=false,                
    showstringspaces=false,
    showtabs=false,                  
    tabsize=2
}
 
\lstset{style=mystyle}
\usepackage{setspace}
\DeclareAcronym{roi}{
short = ROI,
long = region of interest
}
\DeclareAcronym{fig}{
short = Fig.,
long = figure,
class = latex
}
\DeclareAcronym{tab}{
short = Table,
long = table,
class = latex
}
\DeclareAcronym{eq}{
short = Eq.,
long = equation,
class = latex
}
\DeclareAcronym{sec}{
short = Sect.,
long = section,
class = latex
}
\DeclareAcronym{chp}{
short = Chap.,
long = Chapter,
class = latex
}
\DeclareAcronym{fov}{
short = FOV,
long = field of view
}
\DeclareAcronym{bow}{
short = BoW,
long = bag of words
}
\DeclareAcronym{svd}{
short = SVD,
long = singular value decomposition
}
\DeclareAcronym{mse}{
short = MSE,
long = mean squared error
}
\DeclareAcronym{pca}{
short = PCA,
long = principal components analysis
}
\DeclareAcronym{lda}{
short = LDA,
long = linear discriminant analysis
}
\DeclareAcronym{lbp}{
short = LBP,
long = local binary pattern
}
\DeclareAcronym{svm}{
short = SVM,
long = support vector machines
}
\DeclareAcronym{knn}{
short = $k$-NN,
long = $k$-nearest neighbour
}
\DeclareAcronym{nn}{
short = NN,
long = neareast neighbour
}
\DeclareAcronym{hog}{
short = HOG,
long = histogram of oriented gradient
}
\DeclareAcronym{rms}{
  short = RMS,
  long = root mean square
}
\DeclareAcronym{vbl}{
  short = VBL,
  long = visual-based localisation
}
\DeclareAcronym{cnn}{
  short = CNN,
  long = convolutional neural networks
}
\DeclareAcronym{rnn}{
  short = RNN,
  long = recurrent neural networks
}
\DeclareAcronym{dl}{
  short = DL,
  long = deep learning
}
\DeclareAcronym{ml}{
  short = ML,
  long = machine learning
}
\DeclareAcronym{dof}{
  short = DoF,
  long = degrees of freedom
}
\DeclareAcronym{sfm}{
  short = SfM,
  long = structure from motion
}
\DeclareAcronym{dem}{
  short = DEM,
  long = digital elevation model
}
\DeclareAcronym{mtl}{
  short = MTL,
  long = multi-task learning
}
\DeclareAcronym{gmm}{
  short = GMM,
  long = Gaussian mixture model	
}
\DeclareAcronym{lstm}{
  short = LSTM,
  long = long short-term memory
}
\DeclareAcronym{anr}{
  short = ANR,
  long = Agence Nationale de la Recherche
}
\DeclareAcronym{plat}{
  short = pLaTINUM,
  long = Cartographie Long Terme pour la Navigation Urbaine
}
\DeclareAcronym{cbir}{
	short = CBIR,
	long = content-based image retrieval
}


%: ----------------------------------------------------------------------
%:                  TITLE PAGE: name, degree,..
% ----------------------------------------------------------------------
% below is to generate the title page with crest and author name

%if output to PDF then put the following in PDF header
\ifpdf  
    \pdfinfo { /Title  (Vision-based localization with discriminative features from heterogeneous visual data)
               /Creator (TeX)
               /Producer (pdfTeX)
               /Author (Nathan Piasco nathan.piasco@gmail.com)
               /CreationDate (D:201904011200ss)  %format D:YYYYMMDDhhmmss
               /ModDate (D:201904011200ss)
               /Subject (PhD Dissertation of Nathan Piasco)
               /Keywords (Visual localisation, CBIR, multi-modal data) }
    \pdfcatalog { /PageMode (/UseOutlines)
                  /OpenAction (fitbh)  }
\fi


\title{Vision-based localization with discriminative features from heterogeneous visual data}



% ----------------------------------------------------------------------
% The section below defines www links/email for author and institutions
% They will appear on the title page of the PDF and can be clicked
\ifpdf
  \author{\href{mailto:nathan.piasco@u-bourgogne.fr}{Nathan Piasco}}
%  \cityofbirth{born in XYZ} % uncomment this if your university requires this
%  % If city of birth is required, also uncomment 2 sections in PhDthesisPSnPDF
%  % Just search for the "city" and you'll find them.
  % The crest is a graphics file of the logo of your research institution.
  % Place it in ./0_frontmatter/figures and specify the width

%% First university
  \firstlab{\href{http://le2i.cnrs.fr/}{ImViA-VIBOT}}
  \firstlogolab{\includegraphics[width=2cm]{logos/logole2i.eps}}
  \firstuni{\href{http://www.u-bourgogne.fr/}{Universit\'e de Bourgogne}}
  \firstlogouni{\includegraphics[width=2cm]{logos/logo-ub-no-bg.pdf}}

  
%% Second university
  \secondlab{\href{http://recherche.ign.fr/presentation.php}{LaSTIG ACTE, IGN, ENSG}}
  \secondlogolab{\includegraphics[width=2cm]{logos/IGN_log_Q}}

  \seconduni{\href{https://www.udg.edu/}{Universit\'e Paris-Est}}
  \secondlogouni{\includegraphics[width=2cm]{logos/univ-paris-est.png}}

  \supervisora{D\'esir\'e Sidib\'e (ImViA - UBFC)}
  \supervisorb{Val\'erie Gouet-Brunet (LaSTIG - UPE)}
  \supervisorc{C\'edric Demonceaux (ImViA - UBFC)}
  \supervisord{}
% If you are not creating a PDF then use the following. The default is PDF.
\else
  \author{Nathan Piasco}
%  \cityofbirth{born in XYZ}
%% First university
  \firstlab{\href{http://le2i.cnrs.fr/}{ImViA-VIBOT}}
  \firstlogolab{\includegraphics[width=2cm]{logos/logole2i.eps}}
  \firstuni{\href{http://www.u-bourgogne.fr/}{Universit\'e de Bourgogne}}
  \firstlogouni{\includegraphics[width=2cm]{logos/logo-ub-no-bg.pdf}}

  
%% Second university
  \secondlab{\href{http://recherche.ign.fr/presentation.php}{LaSTIG ACTE, IGN, ENSG}}
  \secondlogolab{\includegraphics[width=2cm]{logos/IGN_log_Q}}

  \seconduni{\href{https://www.udg.edu/}{Universit\'e Paris-Est}}
  \secondlogouni{\includegraphics[width=4cm]{logos/univ-paris-est.png}}

  \supervisora{D\'esir\'e Sidib\'e (ImViA - UBFC)}
  \supervisorb{Val\'erie Gouet-Brunet (LaSTIG - UPE)}
  \supervisorc{C\'edric Demonceaux (ImViA - UBFC)}
  \supervisord{}
\fi

%\renewcommand{\submittedtext}{change the default text here if needed}
\degree{Philosophi\ae Doctor (PhD)}
\degreedate{October 2019}


% ----------------------------------------------------------------------
       
% turn of those nasty overfull and underfull hboxes
\hbadness=10000
\hfuzz=50pt


%: --------------------------------------------------------------
%:                  FRONT MATTER: dedications, abstract,..
% --------------------------------------------------------------
\onehalfspacing
%\DeclareInstance{acro-title}{empty}{sectioning}{name-format =}

\AtBeginDocument{%
	\tableofcontents 
}

\AtEndDocument{%
	\bibliographystyle{abbrvnat} % calls style file plainnat.bst
	
	\renewcommand{\bibname}{References} % changes the header; default: Bibliography
	
	%\bibliography{9_backmatter/references} % adjust this to fit your BibTex file
	\bibliography{6_backmatter/mendeley}
}

\begin{document}

%\language{english}

% sets line spacing
%\renewcommand\baselinestretch{1.2}
\baselineskip=18pt plus1pt


%: ----------------------- generate cover page ------------------------

\maketitle  % command to print the title page with above variables

%: ----------------------- cover page back side ------------------------
% Your research institution may require reviewer names, etc.
% This cover back side is required by Dresden Med Fac; uncomment if needed.

\newpage
{\pagestyle{plain}
\vspace{10mm}
Reviewers:

\begin{itemize}
\item[] 
\item[] 
%\item[] Soumya Ghose, Research Associate at Case Western Reserve University
\end{itemize}

\vspace{20mm}
Day of the defense: 

\vspace{20mm}
\hspace{70mm}Signature from head of PhD committee:



%: ----------------------- abstract ------------------------

% Your institution may have specific regulations if you need an abstract and where it is to be placed in the document. The default here is just after title.

%
% Thesis Abstract -----------------------------------------------------


%\begin{abstractslong}    %uncommenting this line, gives a different abstract heading
\begin{abstracts}        %this creates the heading for the abstract page

\end{abstracts}
%\end{abstractlongs}
%-------------------------------------------------------------------------

\begin{abstractFrench}

\end{abstractFrench}


% The original template provides and abstractseparate environment, if your institution requires them to be separate. I think it's easier to print the abstract from the complete thesis by restricting printing to the relevant page.
%\begin{abstractseparate}
%  
% Thesis Abstract -----------------------------------------------------


%\begin{abstractslong}    %uncommenting this line, gives a different abstract heading
\begin{abstracts}        %this creates the heading for the abstract page

\end{abstracts}
%\end{abstractlongs}
%-------------------------------------------------------------------------

\begin{abstractFrench}

\end{abstractFrench}

%\end{abstractseparate}


%: ----------------------- tie in front matter ------------------------

% Thesis Dedictation ---------------------------------------------------

\begin{dedication} %this creates the heading for the dedication page
\begin{flushright}
\textit{.}
\end{flushright}
\end{dedication}

% ----------------------------------------------------------------------

\begin{publication}

\subsubsection*{Peer-Review Journals Papers}

\begin{enumerate}\scriptsize
	\item \textcolor{red}{\textbf{N. Piasco}, D. Sidib\'e, V. Gouet-Brunet and C. Demonceaux. Auxiliary Modality as Robust Training Signal for
Visual Localization in Challenging Conditions.}
	\item \textbf{N. Piasco}, D. Sidib\'e, C. Demonceaux and V. Gouet-Brunet. A Survey on Visual-Based Localization: On the Benefit of Heterogeneous Data. Pattern Recognition, Volume 74, February 2018, pp.90-109.
\end{enumerate}

\subsubsection*{Peer-Review International Conferences}

\begin{enumerate}\scriptsize
	\item \textbf{N. Piasco}, D. Sidib\'e, C. Demonceaux and V. Gouet-Brunet. Perspective-n-Learned-Point: Pose
Estimation from Relative Depth. 2019 British Machine Vision Conference (BMVC), Cardiff, United Kingdom, September 2019.
	\item \textbf{N. Piasco}, D. Sidib\'e, C. Demonceaux and V. Gouet-Brunet. Geometric Camera Pose Refinement with Learned Depth Maps. 2019 IEEE International Conference on Image Processing (ICIP), Taipei, Taiwan, September 2019.
	\item \textbf{N. Piasco}, D. Sidib\'e, V. Gouet-Brunet and C. Demonceaux. Learning Scene Geometry for Visual Localization in Challenging Conditions. 2019 IEEE International Conference of Robotics and Automation (ICRA), Montreal, Canada, May 2019.	
\end{enumerate}

\subsubsection*{Peer-Review National Conferences}

\begin{enumerate}\scriptsize
	\item \textbf{N. Piasco}, D. Sidib\'e, V. Gouet-Brunet and C. Demonceaux. Apprentissage de modalit\'es auxiliaires pour la localisation bas\'ee vision. Reconnaissance des Formes, Image, Apprentissage et Perception (RFIAP), Champs-sur-Marne, France, June 2018
	\item \textbf{N. Piasco}, D. Sidib\'e, V. Gouet-Brunet and C. Demonceaux. Localisation Bas\'ee Vision : de l’h\'et\'erog\'en\'eit\'e des approches et des données. ORASIS - Journ\'ees francophones des jeunes chercheurs en vision par ordinateur (ORASIS), Colleville-sur-Mer, France, June 2017.
\end{enumerate}

\subsubsection*{Thesis}

\begin{enumerate}\scriptsize
\item \textbf{N. Piasco}. Vision-based localization with discriminative features from heterogeneous visual data.
\end{enumerate}

\end{publication}


%: ----------------------- contents ------------------------

\setcounter{secnumdepth}{3} % organisational level that receives a numbers
\setcounter{tocdepth}{3}    % print table of contents for level 3
% levels are: 0 - chapter, 1 - section, 2 - subsection, 3 - subsection
\clearpage

%: ----------------------- list of acronyms/figures/tables ------------------------

%\begin{abbreviations}
\addcontentsline{toc}{chapter}{List of Abbreviations}
\printacronyms[exclude-classes=latex, name=List of Abbreviations, heading=chapter*]
%\end{abbreviations}

\listoffigures	% print list of figures

\listoftables  % print list of tables


% Thesis Acknowledgements ------------------------------------------------


%\begin{acknowledgementslong} %uncommenting this line, gives a different acknowledgements heading
\begin{acknowledgements}      %this creates the heading for the acknowlegments

\end{acknowledgements}
%\end{acknowledgmentslong}

% ------------------------------------------------------------------------




%\tableofcontents            % print the table of contents


% Thesis Abstract -----------------------------------------------------


%\begin{abstractslong}    %uncommenting this line, gives a different abstract heading
\begin{abstracts}        %this creates the heading for the abstract page

\end{abstracts}
%\end{abstractlongs}
%-------------------------------------------------------------------------

\begin{abstractFrench}

\end{abstractFrench}

}


%: --------------------------------------------------------------
%:                  MAIN DOCUMENT SECTION
% --------------------------------------------------------------

% the main text starts here with the introduction, 1st chapter,...
\mainmatter

%\renewcommand{\chaptername}{} % uncomment to print only "1" not "Chapter 1"


%: ----------------------- subdocuments ------------------------
% Parts of the thesis are included below. Rename the files as required.
% But take care that the paths match. You can also change the order of appearance by moving the include commands.

\subfile{1_introduction/chapter_1}	% background information + intro
\subfile{2_visual_based_localization/chapter_2} % SOTA
\subfile{3_side_modality_learning_for_localisation/chapter_3} % CBIR VLB method
\subfile{4_pose_refinement/chapter_4} % Relative Pose refinement
\subfile{5_conclusion/chapter_5} %  Conclusion + perspectives

% --------------------------------------------------------------
%:                  BACK MATTER: appendices, refs,..
% --------------------------------------------------------------

% the back matter: appendix and references close the thesis


%: ----------------------- bibliography ------------------------

% The section below defines how references are listed and formatted
% The default below is 2 columns, small font, complete author names.
% Entries are also linked back to the page number in the text and to external URL if provided in the BibTex file.

% PhDbiblio-url2 = names small caps, title bold & hyperlinked, link to page 
%\begin{multicols}{2} % \begin{multicols}{ # columns}[ header text][ space]
%\begin{tiny} % tiny(5) < scriptsize(7) < footnotesize(8) < small (9)

%\bibliographystyle{Latex/Classes/PhDbiblio-url2} % Title is link if provided

%\end{tiny}
%\end{multicols}

% --------------------------------------------------------------
% Various bibliography styles exit. Replace above style as desired.

% in-text refs: (1) (1; 2)
% ref list: alphabetical; author(s) in small caps; initials last name; page(s)
%\bibliographystyle{Latex/Classes/PhDbiblio-case} % title forced lower case
%\bibliographystyle{Latex/Classes/PhDbiblio-bold} % title as in bibtex but bold
%\bibliographystyle{Latex/Classes/PhDbiblio-url} % bold + www link if provided

%\bibliographystyle{Latex/Classes/jmb} % calls style file jmb.bst
% in-text refs: author (year) without brackets
% ref list: alphabetical; author(s) in normal font; last name, initials; page(s)

%\bibliographystyle{plainnat} % calls style file plainnat.bst
% in-text refs: author (year) without brackets
% (this works with package natbib)


% --------------------------------------------------------------

% according to Dresden med fac summary has to be at the end
%
% Thesis Abstract -----------------------------------------------------


%\begin{abstractslong}    %uncommenting this line, gives a different abstract heading
\begin{abstracts}        %this creates the heading for the abstract page

\end{abstracts}
%\end{abstractlongs}
%-------------------------------------------------------------------------

\begin{abstractFrench}

\end{abstractFrench}


%: Declaration of originality
% \include{9_backmatter/declaration}



\end{document}
